\documentclass[10pt,a4paper]{amsart}
\usepackage[margin=19mm]{geometry}
\usepackage{amsmath,amssymb,mathtools}
\usepackage[scr=rsfs,cal=euler]{mathalfa}
\usepackage{xfrac}
\usepackage[unicode,hidelinks,bookmarks=false]{hyperref}
\newcommand{\ie}{i.e.\ }
\newcommand{\cf}{cf.\ }
\newcommand{\resp}{respectively\ }
\newcommand{\Subsection}{\subsection}
\providecommand{\nobreakdash}{\nobreak\hbox{-}}
\setlength{\emergencystretch}{2em}
\multlinegap0pt
\usepackage{amssymb}
\usepackage{mathtools}
\usepackage[shortlabels]{enumitem}
\newtheorem{proposition}{Proposition}
\DeclareMathOperator{\Ric}{Ric}
\newcommand{\mA}{\mathcal{A}}
\title{intermediate curvature and splitting theorem}
\author{Jingche Chen, Han Hong}
\date{Source: arXiv:2604.26529v1 (CC BY 4.0)}
\begin{document}
\maketitle

\noindent\emph{Source and licence.} intermediate curvature and splitting theorem. Version arXiv:2604.26529v1 (CC BY 4.0), distributed by arXiv under the Creative Commons Attribution 4.0 International licence, http://creativecommons.org/licenses/by/4.0/, as recorded in the arXiv licence metadata for this exact version and shown on the abstract page https://arxiv.org/abs/2604.26529v1. Full licence notice: \texttt{licenses/arxiv-2604.26529-CC-BY-4.0.txt}.\par\smallskip

\noindent\textit{Editorial scope.} Counted unit: the article's proposition C\_m (source0.tex lines 1098--1113). The article's complete original proof of it is delivered with it (source0.tex lines 1114--1149, one \texttt{proof} environment of 36 lines). No mathematics is written, repaired or re-derived: the statement and the proof are the article's own lines, in the article's own notation, and no retained line is substituted or re-flowed.  No further source-local context is needed: every cross-reference of every retained line resolves inside the two delivered spans.  Nothing was omitted from any retained span, so the package carries no omission record.  No retained line contains a citation, so the wrapper carries no bibliography and no bibliography entry.  No retained line contains a figure environment, an image-inclusion command or drawing code, so no image is delivered and the package declares no asset dependency.  Editorial formatting only, in addition: an AMS \texttt{amsart} preamble that restates the article's own declarations and macros used by the retained lines, namely the environments \texttt{proposition} on separate counters, together with the standard packages those lines need.  The retained environments print in this document as Proposition 1 in order of appearance, whereas the article prints the same items with the numbering of its own full document.  The complete native source of the article as distributed in the arXiv e-print for this exact version is bundled unchanged as \texttt{sources/199401/source0.tex}.
\begin{proposition}\label{lift C_m-1 to C_m}
Let $\Sigma$ be an $n$-dimensional Riemannina manifold and let $0\leq k <4$. If there exists a positive smooth function $u$ satisfies $-\frac{4}{4-k}\Delta_{\Sigma}u+C_{m-1}^{\Sigma}u\geq \lambda u$ for $\lambda\geq 0$. Then on the $(n+1)$-dimensional manifold $M=\Sigma\times \mathbb{S}^{1}$ with metric $g_{M}=g_{\Sigma}+u^{4\frac{2-k}{4-k}}d \theta^2$ we have 
\[
-k\Delta_{M}u^{\frac{2}{4-k}}+C_{m}^{M}(e_{1},...,e_{m-1},e_{\theta})u^{\frac{2}{4-k}} \geq \lambda u^{\frac{2}{4-k}}
\]
where by abuse of notation, $u$ is the constant extension  in the $\mathbb{S}^1$-direction, i.e., $u(\cdot,\theta)=u(\cdot)$, $e_\theta$ is the unit vector along  $\mathbb{S}^1$-direction, i.e., $e_\theta=u^{2\frac{k-2}{4-k}}\partial_\theta$ and $e_i, i=1,\cdots,m-1$ are unit tangent vectors of $\Sigma.$

Moreover, $\forall\ \alpha\in \mathbb{R}$, $\Sigma$ is $\mathcal{A}(\Sigma)=\int_{\Sigma}u^{\alpha}$-stable in $M$, i.e. $\forall\ \psi \in C_{0}^{\infty}(\Sigma)$, we have
\begin{equation*}
        \begin{split}
            \frac{d^2}{d t^2}\bigg|_{t=0}\mA(\Sigma_{t})&=\int_{\Sigma}|\nabla^{\Sigma}\psi|^2 u^{\alpha}-\left(|A_{\Sigma}|^2+\Ric_M(\nu,\nu)\right)\psi^2 u^{\alpha}\\
            &+\int_\Sigma \alpha(\alpha-1)u^{\alpha-2}\left(\nabla^M_{\nu}u\right)^2\psi^2 +\int_{\Sigma}\alpha u^{\alpha-1}\left(\Delta_M u-\Delta_{\Sigma} u\right)\psi^2\\
            &\geq 0.
        \end{split}
    \end{equation*}
\end{proposition}

\begin{proof}
    For convenience, we denote $\gamma=\frac{2}{4-k}$ and $\delta=\frac{4-2k}{4-k}$.
    It is not hard to check $\Sigma$ is totally geodesic in $M$ and $\Ric^{M}(e_{\theta},e_{\theta} )=-\frac{\Delta_{\Sigma}(u^{\delta})}{u^{\delta}}$. Thus, we have
    \[
    \begin{aligned}
    &\quad -k\Delta_{M}(u^{\gamma})+C_{m}^{M}(e_{1},...,e_{m-1},e_{\theta})u^{\gamma}\\
    &=-k\big(\Delta_{\Sigma}(u^{\gamma})+\frac{1}{u^{\delta}}\langle \nabla_{\Sigma}(u^{\delta}),\nabla_{\Sigma}(u^{\gamma})\rangle\big) +\big(C_{m-1}^{\Sigma}(e_{1},...,e_{m-1})-\frac{\Delta_{\Sigma}(u^{\delta})}{u^{\delta}}\big)u^{\gamma}   \\
    &\geq -k\big(\Delta_{\Sigma}(u^{\gamma})+\frac{1}{u^{\delta}}\langle \nabla_{\Sigma}(u^{\delta}),\nabla_{\Sigma}(u^{\gamma})\rangle\big) +\big(\frac{4}{4-k}\frac{\Delta_{\Sigma}u}{u}+\lambda-\frac{\Delta_{\Sigma}(u^{\delta})}{u^{\delta}}\big)u^{\gamma}\\
    &=\lambda u^\gamma
    \end{aligned}
    \]
    where in the inequality we have used $-\frac{4}{4-k}\Delta_{\Sigma}u+C_{m-1}^{\Sigma}u\geq \lambda u$.

Now we prove $\Sigma$ is weighted stable. Actually, Since $A_{\Sigma}=0$ and $u_{\nu}=0$, it suffices to show
\begin{equation*}
        \begin{split}
            \int_{\Sigma}|\nabla^{\Sigma}\psi|^2 u^{\alpha}-\Ric_M(e_{\theta},e_{\theta})\psi^2 u^{\alpha} +\int_{\Sigma}\alpha u^{\alpha-1}\left(\Delta_M u-\Delta_{\Sigma} u\right)\psi^2\geq 0.
        \end{split}
    \end{equation*}
Since $\Ric^{M}(e_{\theta},e_{\theta} )=-\frac{\Delta_{\Sigma}(u^{\delta})}{u^{\delta}}$ and $\Delta_{M}u=\Delta_{\Sigma}u+\frac{1}{u^{\delta}}\langle \nabla_{\Sigma}u^{\delta},\nabla_{\Sigma}u\rangle$, we only need to show
\begin{equation*}
        \begin{split}
            \int_{\Sigma}|\nabla^{\Sigma}\psi|^2 u^{\alpha}+\frac{\Delta_{\Sigma}(u^{\delta})}{u^{\delta}}\psi^2 u^{\alpha} +\int_{\Sigma}\alpha u^{\alpha-1}\left(\frac{1}{u^{\delta}}\langle \nabla_{\Sigma}u^{\delta},\nabla_{\Sigma}u\rangle\right)\psi^2\geq 0.
        \end{split}
    \end{equation*}
However, we obtain
\begin{equation*}
        \begin{split}
            &\quad\int_{\Sigma}|\nabla^{\Sigma}\psi|^2 u^{\alpha}+\frac{\Delta_{\Sigma}(u^{\delta})}{u^{\delta}}\psi^2 u^{\alpha} +\int_{\Sigma}\alpha u^{\alpha-1}\left(\frac{1}{u^{\delta}}\langle \nabla_{\Sigma}u^{\delta},\nabla_{\Sigma}u\rangle\right)\psi^2\\
            &=\int_{\Sigma}|\nabla^{\Sigma}\psi|^2 u^{\alpha}+\left(\Delta_{\Sigma}\log u^{\delta}+|\nabla_{\Sigma}\log u^{\delta}|^2\right)\psi^2 u^{\alpha} +\int_{\Sigma}\alpha\delta u^{\alpha-2}|\nabla_{\Sigma}u|^2\psi^2\\
            &=\int_{\Sigma}|\nabla^{\Sigma}\psi|^2 u^{\alpha}+|\nabla_{\Sigma}\log u^{\delta}|^2\psi^2 u^{\alpha} -\langle \nabla_{\Sigma}\log u^{\delta},2\psi\nabla_{\Sigma}\psi\rangle u^{\alpha}\\
            &\geq 0,
        \end{split}
    \end{equation*}
    where the second equation is due to integration by part.
\end{proof}

\end{document}
